% ============================================================================
% Fichier  : partie5/P05_C22_ia_forensique.tex
% Titre    : Intelligence Artificielle et Forensique
% Auteur   : MINKA MI NGUIDJOÏ Thierry Emmanuel
% Version  : 1.0 -- Juin 2026
% Statut   : RÉDIGÉ
% Sources  : Construit ex-nihilo (aucune source projet ne couvre l'IA/ML)
%            Ancré sur le cas MUNGO (deepfakes électoraux, Ch.20) et la
%            méthode ACH/E8-E9 du PIU (Ch.6) appliquée au triage automatisé
% Positionnement : Deuxième chapitre de la Partie V. Deux faces : l'IA comme
%                  OUTIL de l'investigateur (triage, anomalies) et l'IA comme
%                  OBJET d'investigation (deepfakes, contenus synthétiques).
%                  Referme la question laissée ouverte par le cas MUNGO
%                  (Ch.20) sur l'analyse forensique des deepfakes audio.
% ============================================================================

\chapter{Intelligence Artificielle et Forensique}
\label{chap:ia-for}

\epigraph{%
    « L'intelligence artificielle ne ment pas mieux que l'humain ---
    elle ment plus vite, à plus grande échelle, et laisse des traces
    différentes. »%
}{--- Adapté de la pratique forensique contemporaine}

\niveauChapitre{Expert}{%
    Notions de machine learning recommandées mais non indispensables~:
    ce chapitre explique les concepts nécessaires. Lecture du cas
    MUNGO (Ch.~\ref{chap:cas-int}) recommandée.%
}

\begin{boxobjectifs}
\begin{itemize}
    \item Distinguer les deux usages de l'IA en forensique~: l'IA
          comme outil d'investigation et l'IA comme objet d'investigation.
    \item Appliquer des techniques de triage automatisé pour
          prioriser l'analyse de grands volumes de données.
    \item Détecter forensiquement les contenus générés ou modifiés
          par IA~: images, audio (deepfakes), texte.
    \item Comprendre les limites épistémologiques des outils d'IA
          en contexte judiciaire~: biais, explicabilité, robustesse
          au contre-interrogatoire.
    \item Documenter une analyse assistée par IA dans le respect
          de la distinction constatation/analyse (NIST~SP~800-86).
    \item Anticiper les évolutions de l'IA générative et leur
          impact sur l'admissibilité des preuves numériques.
\end{itemize}
\end{boxobjectifs}

\bigskip

% ============================================================================
\section*{Deux visages de l'intelligence artificielle en forensique}
% ============================================================================

Le cas MUNGO (Ch.~\ref{chap:cas-int}) a posé une question restée
ouverte~: comment analyser forensiquement un clip audio deepfake~?
Ce chapitre y répond en élargissant le cadre~: l'intelligence
artificielle se présente sous deux visages distincts pour
l'investigateur numérique.

\begin{enumerate}
    \item \textbf{L'IA comme outil}~: accélérer le triage de
          téraoctets de données, détecter des anomalies invisibles
          à l'\oe il humain, automatiser des tâches répétitives.
    \item \textbf{L'IA comme objet}~: les contenus générés par IA
          (deepfakes, textes synthétiques, images générées) deviennent
          eux-mêmes des preuves ou des éléments à authentifier.
\end{enumerate}

Ces deux visages exigent des compétences différentes mais
complémentaires.

\separateur

% ============================================================================
\section{L'IA comme Outil~: Triage et Détection d'Anomalies}
\label{sec:ia-outil}
% ============================================================================

\subsection{Le problème du volume}

Une investigation moderne implique fréquemment des volumes de données
qu'aucune analyse manuelle exhaustive ne peut traiter dans des délais
raisonnables~: des millions d'emails, des téraoctets de logs réseau,
des dizaines de milliers de photos. Le rôle de l'IA n'est pas de
remplacer l'investigateur mais de \textbf{prioriser} ce qu'il doit
examiner en premier.

\begin{boxdef}
\textbf{Triage assisté par IA}\index{triage assisté par IA}~:
utilisation d'algorithmes de classification ou de détection
d'anomalies pour ordonner un volume de données selon leur pertinence
probable, \emph{sans} se substituer à la décision finale de
l'investigateur humain.
\end{boxdef}

\subsection{Classification d'emails et de documents}

\begin{boxoutils}{Triage d'emails avec classification supervisée}
\begin{lstlisting}[language=bash_forensique]
# Cas type : 50 000 emails saisis dans l'affaire FAKO (Ch.20)
# Objectif : identifier les emails pertinents pour la fraude
# comptable sans les lire un par un

# Approche 1 : classification par mots-cles ponderes (simple,
# explicable -- recommande en premiere passe)
python3 << 'EOF'
import re
from collections import Counter

mots_cles = {
    'virement': 3, 'compte bancaire': 3, 'urgent': 1,
    'confidentiel': 2, 'ne pas divulguer': 3, 'IBAN': 3,
    'mot de passe': 2, 'autorisation': 1
}

def score_email(texte):
    texte_lower = texte.lower()
    return sum(poids for mot, poids in mots_cles.items()
               if mot in texte_lower)

# Trier les emails par score decroissant
# Les 200 emails les plus eleves sont examines en priorite
EOF

# Approche 2 : modele de classification (necessite jeu
# d'entrainement -- moins explicable, a documenter avec precaution)
# scikit-learn, classification Naive Bayes ou SVM
# ATTENTION : un modele "boite noire" pose un probleme
# d'explicabilite en contexte judiciaire (cf. section limites)
\end{lstlisting}
\end{boxoutils}

\begin{boxCRO}
\textbf{Analyse CRO --- Triage IA et risque de biais}

$\mathcal{R}$ (Fiabilité)~: un modèle de triage bien calibré accélère
l'identification de preuves pertinentes, mais peut introduire des
faux négatifs systématiques (documents pertinents non identifiés
car ne correspondant pas aux patterns d'entraînement).

$\mathcal{O}$ (Opposabilité)~: un tri qui \emph{exclut} des documents
de l'examen humain pose un risque majeur~: si la défense découvre
qu'un document à décharge a été écarté par l'algorithme sans examen
humain, c'est l'ensemble du processus qui est compromis.

\cro{$\sim$}{$\uparrow$ vitesse, $\downarrow$ risque faux négatifs}{$\downarrow\downarrow$ si exclusion sans validation humaine}

\textbf{Règle pratique impérative~:} l'IA \textbf{priorise}, elle
ne \textbf{décide jamais} de l'exclusion définitive d'un document
de l'examen. Tout document, même classé comme peu pertinent, doit
rester accessible et son score de classification documenté.
\end{boxCRO}

\subsection{Détection d'anomalies dans les logs}

\begin{boxoutils}{Détection d'anomalies comportementales (UEBA)}
\begin{lstlisting}[language=bash_forensique]
# UEBA : User and Entity Behavior Analytics
# Principe : etablir une baseline du comportement normal d'un
# utilisateur ou systeme, puis detecter les ecarts statistiques

# Exemple simplifie : detection d'horaires de connexion anormaux
python3 << 'EOF'
import pandas as pd
from sklearn.ensemble import IsolationForest

# Chargement des logs de connexion (Event ID 4624, Ch. 14)
df = pd.read_csv('connexions.csv')
# Colonnes : heure_connexion, jour_semaine, duree_session, volume_donnees

# Modele de detection d'anomalies (non supervise)
modele = IsolationForest(contamination=0.05, random_state=42)
df['anomalie'] = modele.fit_predict(
    df[['heure_connexion', 'jour_semaine', 'duree_session']]
)

# Les connexions marquees -1 sont des anomalies statistiques
anomalies = df[df['anomalie'] == -1]
print(f"{len(anomalies)} connexions anormales identifiees sur "
      f"{len(df)} totales")
EOF

# IMPORTANT : une anomalie statistique n'est PAS une preuve de
# malveillance. C'est un signal de priorisation pour examen humain.
# Documenter : "X connexions ont ete identifiees comme statistiquement
# atypiques par le modele Isolation Forest (parametres : contamination=
# 0.05) et ont fait l'objet d'un examen manuel detaille."
\end{lstlisting}
\end{boxoutils}

\subsection{Reconnaissance d'objets dans les preuves visuelles}

\begin{boxoutils}{Détection automatisée dans des collections de médias}
\begin{lstlisting}[language=bash_forensique]
# Cas type : 15 000 photos extraites d'un telephone (Ch. 16)
# Objectif : identifier rapidement les photos contenant des
# documents, des armes, ou des visages specifiques

# Outils disponibles dans Autopsy (Ch. 16) :
# - Module "Object Detection" : categories generiques
#   (personne, vehicule, document, arme)
# - Module "Facial Recognition" : regroupement par visage
#   (PAS d'identification -- juste regroupement de similarite)

# Limite cruciale : la reconnaissance faciale automatisee a un
# taux d'erreur variable selon les populations -- documenter
# systematiquement le taux de faux positifs/negatifs du modele
# utilise et ne JAMAIS presenter une correspondance faciale
# automatisee comme une identification certaine sans validation
# par un examinateur qualifie
\end{lstlisting}
\end{boxoutils}

% ============================================================================
\section{L'IA comme Objet~: Détecter les Contenus Synthétiques}
\label{sec:ia-objet}
% ============================================================================

\subsection{Taxonomie des contenus générés par IA}

\begin{tableaucours}{p{2.8cm}p{3.5cm}p{5.0cm}}{%
    \tableauHeader{Type} &
    \tableauHeader{Technologie typique} &
    \tableauHeader{Contexte forensique camerounais}
}
\textbf{Deepfake vidéo} & GAN, modèles de diffusion vidéo &
    Vidéos compromettantes fabriquées, désinformation politique \\
\textbf{Deepfake audio} & Clonage vocal (modèles \textit{text-to-speech}) &
    Fausses déclarations attribuées (cas MUNGO, Ch.~\ref{chap:cas-int})~;
    arnaques par usurpation vocale \\
\textbf{Image générée} & Diffusion (Midjourney, Stable Diffusion) &
    Faux documents d'identité, fausses preuves visuelles \\
\textbf{Texte généré} & LLM (GPT, Claude, Llama) &
    Faux témoignages écrits, emails de phishing sophistiqués,
    faux articles de presse \\
\end{tableaucours}
\vspace{-0.8em}
\begin{center}{\small\itshape Tableau~22.1 --- Taxonomie des contenus synthétiques}\end{center}

\subsection{Détection de deepfakes audio~: le cas MUNGO résolu}

Reprenons la question laissée ouverte au Ch.~\ref{chap:cas-int}~:
comment analyser forensiquement les 12~clips audio deepfake de
l'affaire MUNGO~?

\begin{boxoutils}{Analyse forensique d'un clip audio suspect}
\begin{lstlisting}[language=bash_forensique]
# Etape 1 : analyse des metadonnees du fichier audio
exiftool clip_suspect.mp3
# Chercher : logiciel d'encodage, absence de metadonnees
# d'enregistrement original (date GPS, modele d'appareil)
# -- un enregistrement authentique sur smartphone a generalement
# des metadonnees riches ; un deepfake genere n'en a souvent aucune

# Etape 2 : analyse spectrale (spectrogramme)
sox clip_suspect.mp3 -n spectrogram -o spectrogramme.png
# Les voix de synthese presentent souvent :
# - une absence de bruit de fond naturel (silence trop "propre")
# - des artefacts de frequence caracteristiques du vocoder utilise
# - une regularite anormale du rythme respiratoire (souvent absent)

# Etape 3 : outils de detection specialises
# (recherche academique active, taux de detection variable)
# - Resemblyzer (analyse d'empreinte vocale)
python3 -c "
from resemblyzer import VoiceEncoder, preprocess_wav
encoder = VoiceEncoder()
wav = preprocess_wav('clip_suspect.wav')
embedding = encoder.embed_utterance(wav)
# Comparer avec un echantillon authentique connu du candidat
"

# - Detecteurs commerciaux : Pindrop, Resemble Detect,
#   Microsoft Video Authenticator (limites : taux de faux
#   positifs/negatifs evoluent avec chaque nouvelle generation
#   de modeles generatifs -- TOUJOURS documenter la version
#   de l'outil et sa date d'evaluation independante)

# Etape 4 : analyse contextuelle (la plus fiable a ce jour)
# - Comparer avec des enregistrements authentiques connus
#   (discours publics du candidat, interviews)
# - Verifier la coherence factuelle du contenu attribue
# - Tracer l'origine de diffusion (premiere apparition,
#   compte source -- cf. Ch. 15, analyse reseau)
\end{lstlisting}
\end{boxoutils}

\begin{boxalert}
\textbf{L'état de l'art~: course technologique sans fin}

La détection de deepfakes est un domaine de recherche en évolution
rapide et \textbf{aucun outil de détection actuel n'offre une
fiabilité absolue}. Chaque amélioration des détecteurs est suivie
d'une amélioration des générateurs (course aux armements
\textit{GAN vs détecteur}). Conséquence pour le rapport forensique~:
\begin{itemize}
    \item Ne jamais affirmer ``ce clip est un deepfake'' avec
          certitude absolue sur la seule base d'un détecteur
          automatisé~;
    \item Toujours croiser l'analyse technique avec l'analyse
          contextuelle (provenance, cohérence factuelle, témoignages)~;
    \item Documenter le taux de faux positifs/négatifs connu
          de l'outil utilisé à la date de l'analyse.
\end{itemize}
\end{boxalert}

\subsection{Détection d'images générées par IA}

\begin{boxoutils}{Indices forensiques d'images générées}
\begin{lstlisting}[language=bash_forensique]
# Indices visuels classiques (modeles de diffusion, 2023-2025) :
# - mains avec un nombre incorrect de doigts (de moins en moins
#   fiable : les modeles recents corrigent ce defaut)
# - incoherences dans les reflets, ombres, perspectives
# - texte illisible ou deforme dans l'image (panneaux, etiquettes)
# - metadonnees EXIF absentes ou incoherentes avec le contexte allegue

exiftool image_suspecte.jpg
# Rechercher specifiquement :
# - Software tag (certains generateurs laissent une signature,
#   ex : "Software: Midjourney" -- de moins en moins frequent
#   car facilement supprime)
# - Absence totale de donnees EXIF (suspect pour une photo
#   pretendue prise avec un smartphone)

# Detection par analyse de frequence (artefacts de generation)
# Outils specialises : 
# - Hugging Face model hub (modeles de detection open source)
# - C2PA (Coalition for Content Provenance and Authenticity)
#   -- standard emergent de signature de provenance des medias,
#   verifie via : c2patool inspect image.jpg

# IMPORTANT : l'absence d'indicateur de generation IA NE PROUVE
# PAS l'authenticite -- elle signifie seulement qu'aucun indicateur
# connu n'a ete detecte par les outils disponibles a la date
# de l'analyse
\end{lstlisting}
\end{boxoutils}

\subsection{Détection de texte généré par IA}

\begin{boxoutils}{Analyse de texte suspecté généré par LLM}
\begin{lstlisting}[language=bash_forensique]
# Contexte : un temoignage ecrit, un email, ou un document
# est suspecte d'avoir ete redige par un LLM plutot que par
# son auteur pretendu

# Indices linguistiques (faible fiabilite individuelle, a
# combiner) :
# - structures syntaxiques tres regulieres
# - absence de fautes typiques de l'auteur pretendu
#   (si l'auteur a un niveau de francais connu, par exemple)
# - vocabulaire incoherent avec le profil socio-educatif declare
# - "perplexite" textuelle anormalement basse (mesure statistique
#   de la previsibilite du texte selon un modele de langage)

# Outils de detection (fiabilite limitee et contestee
# scientifiquement -- a utiliser avec EXTREME prudence en
# contexte judiciaire) :
# - GPTZero, Originality.ai, etc.
# ATTENTION : ces outils ont des taux de faux positifs
# significatifs, particulierement sur les textes d'auteurs
# non-natifs du francais/anglais -- RISQUE DE BIAIS DISCRIMINATOIRE
# a documenter explicitement si ces outils sont utilises

# Approche forensique privilegiee : analyse de la chaine de
# production du document (metadonnees du fichier, historique
# de versions, horodatage de creation vs publication) plutot
# que l'analyse stylometrique seule
\end{lstlisting}
\end{boxoutils}

% ============================================================================
\section{Limites Épistémologiques de l'IA en Contexte Judiciaire}
\label{sec:limites-ia}
% ============================================================================

\subsection{Le problème de l'explicabilité}

\begin{boxdef}
\textbf{Explicabilité (XAI)}\index{explicabilité}~: capacité d'un
système d'IA à justifier sa décision dans des termes compréhensibles
par un humain. Les modèles de \textit{deep learning} complexes
(réseaux de neurones profonds) sont notoirement peu explicables
--- ils produisent un résultat sans qu'on puisse facilement tracer
le raisonnement qui y a mené.
\end{boxdef}

\begin{boxjuris}
\textbf{L'explicabilité et le standard FRE~Rule~702 / témoignage expert}

Le Chapitre~\ref{chap:leg-mond} a présenté les \textit{Federal Rules
of Evidence}. La Rule~702 (méthodologie de l'expert) exige que les
méthodes employées soient \textbf{fiables et vérifiables}. Un modèle
d'IA ``boîte noire'' dont même son concepteur ne peut expliquer le
raisonnement pose un problème direct~: comment l'expert peut-il
attester de la fiabilité d'une méthode qu'il ne peut pas expliquer~?

\textbf{Conséquence pratique~:} privilégier, chaque fois que possible,
des méthodes explicables (classification par règles, scores pondérés
documentés) pour les conclusions qui seront présentées au tribunal.
Réserver les modèles complexes (deep learning) au triage préliminaire,
suivi systématique d'une validation humaine explicable.
\end{boxjuris}

\subsection{Biais algorithmique et discrimination}

\begin{boxalert}
\textbf{Le risque de biais dans les outils forensiques IA}

Les modèles d'IA héritent des biais présents dans leurs données
d'entraînement. En contexte forensique, ce risque est documenté
pour~:
\begin{itemize}
    \item \textbf{Reconnaissance faciale}~: taux d'erreur
          significativement plus élevés pour certaines catégories
          démographiques selon de nombreuses études indépendantes
          (NIST, MIT Media Lab)~;
    \item \textbf{Détection de texte généré par IA}~: biais
          documentés contre les locuteurs non-natifs de la
          langue d'analyse~;
    \item \textbf{Reconnaissance vocale}~: performance variable
          selon les accents régionaux.
\end{itemize}
\textbf{Implication pour le Cameroun~:} les outils d'IA forensique
sont majoritairement entraînés sur des corpus occidentaux.
Leur application à des sujets, des voix ou des textes camerounais
(accents locaux, mélange français/anglais/langues locales)
\textbf{peut présenter des taux d'erreur non documentés}.
Exiger systématiquement, avant utilisation en contexte judiciaire,
une évaluation indépendante du taux d'erreur de l'outil sur des
échantillons représentatifs de la population concernée.
\end{boxalert}

\subsection{Robustesse au contre-interrogatoire}

\begin{boxscenario}{Contre-interrogatoire d'une conclusion assistée par IA}
\textbf{Contexte~:} dans l'affaire MUNGO, l'expert a utilisé un outil
de détection de deepfake (taux de détection annoncé par l'éditeur~:
94\%) pour conclure que 8~des~12~clips audio sont des deepfakes.

\textbf{Q~:} ``Quel est le taux de faux positifs de cet outil~?''\\
\textbf{R~:} ``L'éditeur annonce un taux de faux positifs de~3\%
sur son jeu de test publié. Je n'ai pas de donnée indépendante
sur la performance spécifique pour des voix en français camerounais.
J'ai en conséquence croisé cette analyse avec une comparaison
contextuelle [détails de l'analyse contextuelle].''

\textbf{Q~:} ``Pouvez-vous expliquer pourquoi l'algorithme a classé
ce clip spécifique comme deepfake~?''\\
\textbf{R~:} ``L'outil utilisé fournit un score de confiance global
(87\% dans ce cas) mais ne détaille pas les caractéristiques
spécifiques ayant motivé ce score, ce qui constitue une limite que
je documente dans mon rapport. Pour cette raison, ma conclusion
ne repose pas uniquement sur ce score mais sur la convergence avec
[analyse spectrale manuelle, incohérences contextuelles].''

\textbf{Constat~:} la deuxième réponse illustre la bonne pratique~:
l'expert reconnaît la limite de l'outil ET montre que sa conclusion
ne dépend pas uniquement de cette boîte noire.
\end{boxscenario}

% ============================================================================
\section{Documenter une Analyse Assistée par IA dans le Rapport}
\label{sec:documentation-ia}
% ============================================================================

\begin{lstlisting}[language=bash_forensique,
    caption={Modèle de documentation d'une analyse assistée par IA}]
3.X MÉTHODOLOGIE -- OUTIL DE DÉTECTION ASSISTÉE PAR IA

Outil utilisé        : [Nom, version, éditeur]
Date d'utilisation    : [date]
Type de modèle        : [ex : réseau de neurones convolutif
                          entraîné pour la détection de deepfakes
                          audio]
Performance annoncée  : [taux de détection / faux positifs selon
                          la documentation de l'éditeur, avec
                          référence à la publication ou au rapport
                          technique correspondant]
Validation indépendante : [Oui/Non -- si Non, le mentionner
                          explicitement comme une limite]
Limite identifiée      : [absence de données de performance
                          spécifiques au contexte linguistique/
                          démographique camerounais, le cas échéant]

CONSTATATION (résultat brut de l'outil) :
"L'outil [nom] a attribué un score de confiance de 87% à
l'hypothèse que le clip audio [référence] est généré par
synthèse vocale."
[Ceci EST une constatation : c'est un fait vérifiable -- l'outil
a produit ce score. Ce N'EST PAS une preuve que le clip est
effectivement un deepfake.]

ANALYSE (interprétation avec niveau de confiance) :
"Ce résultat, combiné avec [analyse spectrale manuelle révélant
X, absence de métadonnées d'enregistrement original, incohérence
contextuelle Y], converge vers la conclusion que ce clip audio
est très probablement un contenu synthétique, avec un niveau
de confiance modéré à élevé. Cette conclusion ne repose pas
exclusivement sur le score de l'outil automatisé."
\end{lstlisting}

\separateur

\begin{boxresume}
\textbf{Ce qu'il faut retenir de ce chapitre~:}
\begin{itemize}
    \item \textbf{Deux visages de l'IA}~: outil (triage, détection
          d'anomalies, accélération) et objet (deepfakes, contenus
          synthétiques à authentifier). Compétences différentes,
          rigueur identique.

    \item \textbf{Règle d'or du triage}~: l'IA \textbf{priorise},
          elle ne \textbf{décide jamais} de l'exclusion définitive
          d'un document. Tout document doit rester accessible
          et son score documenté.

    \item \textbf{Détection de deepfakes}~: croiser systématiquement
          analyse technique (métadonnées, spectrogramme, outils
          spécialisés) et analyse contextuelle (provenance,
          cohérence factuelle). Aucun outil n'offre une fiabilité
          absolue~--- course technologique permanente.

    \item \textbf{Explicabilité}~: privilégier les méthodes
          explicables pour les conclusions présentées au tribunal.
          Les modèles ``boîte noire'' posent un problème direct
          au regard du standard FRE~Rule~702.

    \item \textbf{Biais algorithmique~:} les outils d'IA forensique
          sont majoritairement entraînés sur des corpus occidentaux.
          Exiger une évaluation indépendante de leur performance
          sur des échantillons représentatifs avant usage judiciaire
          au Cameroun.

    \item \textbf{Documentation}~: le résultat brut d'un outil IA
          est une \textbf{constatation} (``l'outil a produit ce
          score'')~; son interprétation est une \textbf{analyse}
          (avec niveau de confiance et hypothèses alternatives).
          Ne jamais présenter un score automatisé comme une
          conclusion en soi.

    \item \textbf{Contre-interrogatoire}~: anticiper systématiquement
          les questions sur le taux de faux positifs, l'explicabilité
          et la validation indépendante de tout outil d'IA utilisé.
\end{itemize}
\end{boxresume}

\bigskip

\begin{boxexo}[title={\ding{45}\quad Exercice~22.1 --- Triage vs décision \niveaubase}]
Pour chacune des situations suivantes, déterminez si l'usage de
l'IA décrit respecte la règle ``l'IA priorise, ne décide jamais''~:
\begin{enumerate}[(a)]
    \item Un outil de classification trie 50~000~emails par score
          de pertinence~; l'investigateur examine les 500~premiers
          puis conclut son rapport sans mentionner les~49~500~autres.
    \item Un outil UEBA signale 12~connexions comme anormales sur
          5~000~; l'investigateur documente les 12~connexions
          \emph{et} précise dans son rapport les critères et
          paramètres du modèle utilisé.
    \item Un détecteur de deepfake conclut ``deepfake~: 91\%
          de confiance''~; le rapport cite uniquement ce score
          comme preuve de falsification.
\end{enumerate}
Pour chaque cas, identifiez le risque CRO associé.
\end{boxexo}

\begin{boxexo}[title={\ding{45}\quad Exercice~22.2 --- Analyse de deepfake \niveaumoyen}]
Vous recevez un clip audio de 45~secondes dans lequel une voix
ressemblant à celle du Directeur Financier d'une entreprise
camerounaise autorise un virement de 80~millions~FCFA par téléphone.
\begin{enumerate}[(a)]
    \item Listez les 4~étapes d'analyse forensique que vous
          appliqueriez (référez-vous à la structure de ce chapitre).
    \item Quelles métadonnées chercheriez-vous en premier~? Que
          signifierait leur absence totale~?
    \item Si l'outil de détection automatisé donne un score
          ambigu (52\% de confiance ``deepfake''), quelle est votre
          démarche~? Quelles autres sources mobilisez-vous
          (pensez aux Ch.~15~et~17)~?
    \item Rédigez la section ``Constatation'' et ``Analyse''
          correspondant à votre conclusion, en respectant
          la structure du Ch.~\ref{chap:rapport}.
\end{enumerate}
\end{boxexo}

\begin{boxexo}[title={\ding{45}\quad Exercice~22.3 --- IA et biais en contexte camerounais \niveauexpert}]
L'ANTIC envisage de déployer un outil commercial de reconnaissance
vocale pour automatiser le traitement des plaintes orales déposées
par téléphone, en complément du travail des OPJ.
\begin{enumerate}[(a)]
    \item Quelles questions poseriez-vous à l'éditeur de l'outil
          avant son déploiement, en lien avec les biais algorithmiques
          présentés dans ce chapitre~?
    \item Le Cameroun est un pays multilingue (français, anglais,
          plus de 250~langues locales, avec des accents régionaux
          marqués). Quel risque spécifique cela pose-t-il pour
          un outil entraîné sur des corpus majoritairement
          occidentaux~?
    \item Si cet outil devait être utilisé pour trier les plaintes
          par ordre de priorité, quelles garanties proposeriez-vous
          pour respecter le principe ``l'IA priorise, ne décide
          jamais''~?
    \item Rédigez une recommandation de politique d'usage (5~points
          maximum) pour l'adoption responsable de cet outil par
          l'ANTIC, en vous appuyant sur les principes
          d'explicabilité et de validation indépendante.
\end{enumerate}
\end{boxexo}

\bigskip

\begin{boxinfo}
\textbf{Transition.} L'intelligence artificielle transforme à la
fois les capacités de l'investigateur et la nature des preuves
à analyser. Cette double transformation se prolonge dans le défi
suivant~: l'arrivée de l'informatique quantique, qui menace
à terme l'intégrité cryptographique des preuves numériques
signées aujourd'hui. C'est l'objet du chapitre suivant.
\end{boxinfo}